\section{Implementation and Results}

The network transcription uses analytic derivatives and accepts time-varying
load, PV, and price data directly at each stage. Compact problems retain
FilterDDP's dense null-space backward pass. For feeder models, sparse
Jacobians and local Hessian blocks are instead assembled into the saddle-point
system
\begin{equation}
\begin{bmatrix}H&A^{\mathsf T}\\A&0\end{bmatrix}
\begin{bmatrix}\Delta u\\\Delta\lambda\end{bmatrix}=b,
\label{eq:sparse_kkt}
\end{equation}
which is solved by sparse LU without explicitly constructing a dense QR basis
or reduced Hessian. The dense route has less overhead for small problems; the
sparse route is required when those dense intermediate matrices no longer fit
practical memory. Reported FilterDDP times include its optimization iterations,
not model construction. Centralized references use the exact native
second-order cone equivalent of (\ref{eq:net_soc}) in Gurobi and were checked by
an independent constraint validator.

\subsection{IEEE123C}

The loss-aware model was first tested on IEEE123C (128 buses, 127 branches,
51 batteries, and 51 PV units), with per-stage dimensions
$(n_x,n_u,n_c)=(51,791,562)$. FilterDDP returned normally for
$T\in\{3,6,12,24,48,96\}$; its largest equality residual was
$8.9\times10^{-8}$ and its largest absolute objective gap from the centralized
reference was $5.5\times10^{-4}$. Runtime increased from 74.9~s at $T=3$ to
2021.4~s at $T=96$ (48 and 121 iterations, respectively). The corresponding
centralized solve required 12.57~s at $T=96$. This establishes functional
horizon scaling, but not competitive runtime.

\subsection{Corrected IEEE2522C}

The feeder previously labelled ``IEEE2552C'' was renamed IEEE2522C. Its
one-phase aggregate conversion had also combined aggregate three-phase powers
with a phase-voltage base, making per-unit line impedances nine times too
large. Only the electrical base was corrected, to 7.2~kV on 1~MVA;
physical impedances, powers, DER ratings, voltage limits, losses, and the SOC
model were unchanged. With this correction, cold-start centralized models
were globally optimal and independently validated through $T=96$.

Here $(n_x,n_u,n_c)=(250,13358,10337)$ per stage. The original dense
backward pass could not complete one iteration because it attempted a dense
$13358\times13358$ QR basis and a reduced Hessian of order 3021. The sparse
KKT route removes that storage bottleneck and gives the verified results in
Table~\ref{tab:ieee2522_filterddp}. All runs used full line-loss balances and
the exact SOC constraint, without warm starts.

\begin{table}[H]
\centering
\caption{Sparse FilterDDP on Corrected IEEE2522C}
\label{tab:ieee2522_filterddp}
\setlength{\tabcolsep}{3.5pt}
\begin{tabular}{c r r r r}
\toprule
$T$ & Iter. & Time (s) & $|J-J^\star|$ & Eq. residual \\
\midrule
3  & 56 & 233.098  & $7.34\times10^{-4}$ & $5.63\times10^{-10}$ \\
6  & 82 & 836.933  & $1.30\times10^{-3}$ & $1.75\times10^{-9}$ \\
12 & 79 & 1902.103 & $9.21\times10^{-5}$ & $2.29\times10^{-8}$ \\
24 & 84 & 3068.228 & $1.22\times10^{-3}$ & $3.61\times10^{-8}$ \\
\bottomrule
\end{tabular}
\end{table}

Thus the method is numerically and memory feasible at a 24-period horizon on
this feeder, but remains slow: the $T=24$ run took 51.1~min, versus 27.2~s of
Gurobi solver time for the centralized reference. Reusing UMFPACK's symbolic
factorization did not help: the identical $T=3$ solution slowed from 233.1 to
295.3~s. Better scaling therefore requires a more suitable numerical
factorization or exploitation of feeder structure, rather than symbolic-cache
reuse alone.

\subsection{Large10kC Boundary}

The synthetic large10kC feeder also required a data correction before any OPF
comparison was meaningful. Its aggregate voltage base was restored to
12.47~kV, and uniform synthetic ordinary-line sections were shortened from
$(r,x)=(0.07,0.01)~\Omega$ to $(0.0175,0.0025)~\Omega$; area-boundary
sections use $(0.000025,0.000025)~\Omega$. Powers, DER capacities, voltage
limits, loss terms, and SOC constraints were retained. Cold-start Gurobi
solutions then validated at $T=1$ and $T=3$; the latter has 133035 variables,
objective 2985281.3189, 27 barrier iterations, and 6.304~s solver time.

For FilterDDP at $T=3$, the per-stage dimensions are
$(n_x,n_u,n_c)=(1020,54665,42303)$. Blocking the multiple right-hand-side KKT
solve bounded memory and avoided the earlier immediate allocation failure.
However, the run produced no terminal iteration report after two hours and
was stopped manually. Consequently, no FilterDDP convergence or timing claim
is made for large10kC. The experiment instead identifies sparse numerical
factorization and feedback-solve cost as the present scalability boundary.
